\documentclass[12pt, a4paper]{article}

% ── Packages ──────────────────────────────────────────────────────────────────
\usepackage[utf8]{inputenc}
\usepackage[T1]{fontenc}
\usepackage{amsmath}
\usepackage{graphicx}
\usepackage{hyperref}
\usepackage{booktabs}       % professional-quality tables
\usepackage{array}
\usepackage{geometry}
\usepackage{cite}
\usepackage{url}
\usepackage{listings}       % code listings if needed
\usepackage{xcolor}
\usepackage{lineno}         % line numbers for review
\usepackage{setspace}       % double spacing for review

\geometry{margin=2.5cm}
\doublespacing
\linenumbers

% ── Title block ───────────────────────────────────────────────────────────────
\title{
    \textbf{BoneLogic: A Domain-Specialised Language Reasoning Model \\
    for Bone Biomechanics Knowledge Retrieval}
}

\author{
    Khislatjon Valijonov$^{1}$ \\
    \\
    \small $^{1}$ % TODO: Add affiliation
}

\date{\today}

% ══════════════════════════════════════════════════════════════════════════════
\begin{document}
\maketitle

% ── Abstract ──────────────────────────────────────────────────────────────────
\begin{abstract}
The volume of published bone science literature has grown substantially over the past two decades,
making it increasingly difficult for researchers to retrieve and synthesise domain-specific knowledge
efficiently. General-purpose large language models (LLMs), while capable across many domains, lack
the specialised reasoning required to answer precise questions in bone biomechanics, morphology,
and pathology. We present \textbf{BoneLogic}, a domain-specialised language reasoning model (LRM)
for bone science. BoneLogic combines a curated corpus of \textbf{55,277 peer-reviewed papers}
(1900--2025) collected from the Semantic Scholar academic graph with a two-layer architecture:
a perception layer (LLM + vision-language model) for text and image understanding, and a reasoning
layer (LRM) for hypothesis generation and structured inference. This paper describes the data
collection methodology underpinning BoneLogic, including a novel three-tier PDF URL resolution
pipeline that significantly increases the yield of downloadable full-text papers beyond what
direct API access provides. The resulting corpus represents the largest curated open-access
dataset for bone science to date.

\bigskip
\noindent\textbf{Keywords:} bone biomechanics, domain-specialised language model, knowledge
retrieval, retrieval-augmented generation, academic corpus, Semantic Scholar
\end{abstract}

\newpage

% ══════════════════════════════════════════════════════════════════════════════
\section{Introduction}
\label{sec:intro}

Bone science sits at the intersection of biology, mechanics, materials science, and clinical medicine.
Research in this domain spans bone morphology, structure-function relationships,
biomechanical properties, pathological conditions (e.g.\ osteoporosis, fracture),
medical imaging, biomaterials, and computational simulation. The breadth of this field
means that relevant knowledge is distributed across thousands of journals and tens of
thousands of publications, making comprehensive literature synthesis a significant challenge.

Large language models (LLMs) have demonstrated remarkable capabilities in general-purpose
question answering and text generation \cite{TODO}. However, their application to highly
specialised scientific domains remains limited by two factors. First, general training
corpora under-represent niche domains such as bone science relative to more extensively
documented fields. Second, standard LLMs retrieve knowledge from parametric memory,
which cannot be updated without retraining and provides no grounding to specific
peer-reviewed evidence.

Retrieval-Augmented Generation (RAG) \cite{TODO} addresses the second limitation by
grounding model responses in a retrieved document corpus at inference time. However,
the quality of a RAG system is fundamentally bounded by the quality and coverage of
the underlying corpus. Existing academic search tools are not curated for bone science
and return results mixed with unrelated biomedical literature.

We present \textbf{BoneLogic}, a domain-specialised reasoning system for bone science.
The primary contributions of this paper are:

\begin{itemize}
    \item A \textbf{curated corpus} of 55,277 peer-reviewed bone science papers spanning
          1900--2025, assembled using a structured keyword taxonomy of 133 search terms
          across 17 domain groups.
    \item A \textbf{three-tier PDF URL resolution pipeline} that substantially increases
          full-text PDF yield beyond what Semantic Scholar's API directly provides.
    \item A \textbf{two-layer architecture} (LLM + VLM perception layer; LRM reasoning layer)
          designed for structured scientific inference and hypothesis generation in bone science.
\end{itemize}

% ══════════════════════════════════════════════════════════════════════════════
\section{Related Work}
\label{sec:related}

% TODO: Expand this section

\subsection{Domain-Specialised Language Models}
Several domain-specialised LLMs have been developed for medicine and biology,
including BioBERT \cite{TODO}, PubMedBERT \cite{TODO}, and Med-PaLM \cite{TODO}.
These models are pre-trained or fine-tuned on biomedical literature but are not
specialised for bone science, and focus primarily on clinical text rather than
biomechanical or materials science literature.

\subsection{Academic Knowledge Retrieval}
Semantic Scholar \cite{semanticscholar} provides the largest publicly accessible
index of academic papers, covering over 200 million publications across all
scientific disciplines. PubMed \cite{pubmed} is the standard for biomedical
literature but its coverage of biomechanics and materials science is limited.
Neither system provides domain-filtered retrieval tailored to bone science.

\subsection{Retrieval-Augmented Generation}
RAG systems augment LLMs with a retrieval step, allowing responses to be grounded
in specific documents rather than parametric memory alone \cite{TODO}.
This approach is particularly valuable in scientific domains where factual accuracy
and citation of evidence are essential.

% ══════════════════════════════════════════════════════════════════════════════
\section{Data Collection}
\label{sec:data}

\subsection{Source Selection: Semantic Scholar}

We selected Semantic Scholar \cite{semanticscholar} as the primary data source for
several reasons. First, it maintains the largest publicly accessible database of
academic papers, indexing over 200 million publications across all scientific
fields \cite{semanticscholar}. Second, it provides a free, documented REST API
with structured metadata including title, abstract, authors, year, venue, citation
count, external identifiers (DOI, PubMed ID, arXiv ID), and open-access PDF URLs.
Third, its coverage of biomechanics, materials science, and clinical medicine ---
the three principal disciplines of bone science --- is substantially broader than
domain-specific databases such as PubMed.

Alternative sources were considered but rejected. PubMed's E-utilities API provides
XML full-text access for a subset of papers, but coverage is biased toward clinical
medicine and excludes most biomechanics and materials science literature.
Web of Science and Scopus do not offer free programmatic access.

\subsection{Why PDF Extraction Over API Full Text}

Academic APIs, including the Semantic Scholar API, return only titles and abstracts
in structured form --- not full paper body text. PubMed Central (PMC) provides
XML full text for a subset of open-access papers, but this XML format is
inconsistent across publishers and difficult to parse reliably. We chose instead
to download open-access PDFs and extract full text from the PDF directly.
This approach offers several advantages:

\begin{itemize}
    \item \textbf{Consistent format}: PDF is the universal distribution format
          for academic papers regardless of publisher.
    \item \textbf{Complete content}: PDFs contain figures, tables, equations,
          and supplementary material that XML exports often omit or misformat.
    \item \textbf{Broader coverage}: Open-access PDFs are available from a wider
          range of publishers than PMC XML.
\end{itemize}

\subsection{Keyword Taxonomy}

We constructed a keyword taxonomy of \textbf{133 search queries} organised into
\textbf{17 thematic groups}, covering the full scope of bone science research
(Table~\ref{tab:keywords}).

\begin{table}[h]
\centering
\caption{Keyword taxonomy used for corpus construction. Each group targets a
distinct sub-domain of bone science.}
\label{tab:keywords}
\begin{tabular}{ll}
\toprule
\textbf{Group} & \textbf{Example queries} \\
\midrule
Bone types — skull       & ``cranial bone'', ``skull morphology'' \\
Bone types — spine       & ``vertebral bone'', ``spinal biomechanics'' \\
Bone types — upper limb  & ``humerus'', ``radius bone mechanics'' \\
Bone types — lower limb  & ``femur'', ``tibia stress fracture'' \\
Morphology               & ``bone microstructure'', ``trabecular architecture'' \\
Structure-function       & ``bone anisotropy'', ``cortical bone properties'' \\
Mechanics                & ``bone fracture toughness'', ``bone fatigue'' \\
Pathology                & ``osteoporosis'', ``osteogenesis imperfecta'' \\
Imaging                  & ``bone CT imaging'', ``DXA bone density'' \\
Biomaterials             & ``bone scaffold'', ``hydroxyapatite'' \\
Simulation               & ``finite element bone'', ``bone remodelling simulation'' \\
Cell biology             & ``osteoblast'', ``osteoclast bone resorption'' \\
\multicolumn{2}{c}{\ldots and 5 further groups} \\
\bottomrule
\end{tabular}
\end{table}

Search queries were submitted to the Semantic Scholar API with a rate limit of
one request per second, as required by the API terms of service.
Results were paginated at 100 papers per request.
Deduplication was enforced at the database level using the Semantic Scholar
paper identifier as a primary key, ensuring that papers matching multiple keywords
are stored exactly once.

\subsection{Corpus Statistics}

The complete metadata collection run yielded the statistics shown in
Table~\ref{tab:corpus}.

\begin{table}[h]
\centering
\caption{BoneLogic corpus statistics after Phase 1 data collection.}
\label{tab:corpus}
\begin{tabular}{ll}
\toprule
\textbf{Metric} & \textbf{Value} \\
\midrule
Total papers collected         & 55,277 \\
Unique (no duplicates)         & 100\% (enforced by primary key) \\
Papers with open-access URL    & 21,571 (39\%) \\
PDFs successfully downloaded   & 6,125 \\
Year range                     & 1900 -- 2025 \\
Search keywords                & 133 across 17 topic groups \\
Top journal                    & \textit{Bone} (2,452 papers) \\
\bottomrule
\end{tabular}
\end{table}

The gap between papers with open-access URLs (21,571) and successfully downloaded
PDFs (6,125) reflects papers where the URL exists in the Semantic Scholar database
but the full text requires institutional access at the publisher's website.
This is a known limitation of open-access metadata --- a URL may be classified
as open-access at the time of indexing but subsequently paywalled,
or may require login despite being listed as freely available.

\subsection{PDF URL Resolution Pipeline}

Semantic Scholar's \texttt{openAccessPdf} field contains a heterogeneous mix of
URL types: direct PDF links, DOI redirects, publisher landing pages, and HTML
viewer pages. Attempting to download these URLs directly yields a high proportion
of HTML error pages and login walls rather than actual PDF content.

To address this, we developed a \textbf{three-tier URL resolution pipeline}:

\begin{enumerate}
    \item \textbf{Tier 1 --- Direct PDF detection}: URLs already pointing directly
          to a PDF file (identified by \texttt{.pdf} extension or known
          direct-download path patterns) are used as-is.
    \item \textbf{Tier 2 --- Publisher-specific URL rewriting}: Seven major
          publishers use predictable URL structures that can be algorithmically
          transformed to direct PDF links. Supported publishers include
          Europe PMC, PubMed Central, PLOS, Frontiers, MDPI, Wiley, and Springer.
    \item \textbf{Tier 3 --- Unpaywall API lookup}: For URLs that remain
          unresolved after Tiers 1 and 2, the paper's DOI is submitted to
          the Unpaywall API \cite{unpaywall}, a free open-access PDF index
          covering over 50 million papers.
\end{enumerate}

All downloaded files are validated by checking the first four bytes for the
PDF magic number (\texttt{\%PDF}). Files that fail this check (i.e.\ HTML
pages silently served with a 200 OK status) are discarded.
This pipeline substantially increases the yield of valid PDF downloads
compared to naive direct-download attempts.

% ══════════════════════════════════════════════════════════════════════════════
\section{System Architecture}
\label{sec:architecture}

% TODO: Expand with Phase 2 (RAG) and Phase 3 (LRM) details as built

BoneLogic is designed as a two-layer system:

\begin{itemize}
    \item \textbf{Layer 1 — Perception}: A large language model (LLM) for
          text understanding combined with a vision-language model (VLM)
          for processing bone imaging data (X-ray, CT, MRI, histology).
    \item \textbf{Layer 2 — Reasoning}: A language reasoning model (LRM)
          that operates over retrieved evidence to generate structured
          inferences and testable hypotheses in bone science.
\end{itemize}

The knowledge base underpinning Layer 1 is constructed from the corpus
described in Section~\ref{sec:data}, retrieved via a
Retrieval-Augmented Generation (RAG) pipeline (to be described in full
as Phase 2 is completed).

% ══════════════════════════════════════════════════════════════════════════════
\section{Discussion}
\label{sec:discussion}

% TODO: Fill in as experiments are conducted

\subsection{Limitations}
The current corpus is limited to papers accessible through the Semantic Scholar
API and open-access PDFs. Papers behind institutional paywalls, while present
in the metadata, could not be included in the full-text corpus.
This introduces a potential bias toward open-access journals and preprints.

\subsection{Future Work}
Subsequent phases of BoneLogic will address:
\begin{itemize}
    \item Full-text extraction from downloaded PDFs (PyMuPDF-based chunking pipeline)
    \item Dense embedding of paper chunks into a vector store for RAG retrieval
    \item Integration of bone science textbooks as a complementary knowledge source
    \item Fine-tuning or prompting of the LRM layer for hypothesis generation
    \item A feedback loop for iterative model improvement based on expert validation
\end{itemize}

% ══════════════════════════════════════════════════════════════════════════════
\section{Conclusion}
\label{sec:conclusion}

% TODO: Write once more sections are complete

We have presented the data collection methodology for BoneLogic, a domain-specialised
reasoning system for bone science. The resulting corpus of 55,277 papers, assembled
through a structured keyword taxonomy and a novel three-tier PDF resolution pipeline,
represents a foundation for domain-specialised retrieval and reasoning in bone biomechanics.
Subsequent phases will extend this work to full-text RAG retrieval and structured
hypothesis generation.

% ══════════════════════════════════════════════════════════════════════════════
\section*{Acknowledgements}
% TODO

% ══════════════════════════════════════════════════════════════════════════════
\bibliographystyle{unsrt}
\bibliography{references}

\end{document}
